% Generated draft: requires mathematical and source review.
\documentclass[11pt]{article}
\usepackage[margin=25mm]{geometry}
\usepackage{fontspec}
\setmainfont{Latin Modern Roman}
\usepackage{amsmath,amssymb,mathtools,unicode-math}
\setmathfont{Latin Modern Math}
\usepackage{enumitem,needspace,xurl,hyperref}
\setlength{\parindent}{0pt}
\setlength{\parskip}{5pt}
\setlength{\emergencystretch}{3em}
\newcommand{\N}{\mathbb{N}}
\newcommand{\Z}{\mathbb{Z}}
\newcommand{\Q}{\mathbb{Q}}
\newcommand{\R}{\mathbb{R}}
\newcommand{\C}{\mathbb{C}}
\newcommand{\F}{\mathbb{F}}
\newcommand{\A}{\mathbb{A}}
\newcommand{\PP}{\mathbb{P}}
\newcommand{\E}{\mathbb{E}}
\newcommand{\eps}{\varepsilon}
\newcommand{\dd}{\,\mathrm{d}}
\begin{document}
\section*{UnsolvedMath ID: 2100409}
\textbf{Source problem code:} catalogue entry
\subsection{Definitions and mathematical statement}
V. Retakh's Problem 4.9 has two tasks:
\begin{enumerate}
\item Establish complete integrability of the source's noncommutative version of the leapfrog map.
\item Define noncommutative versions of the pentagram map and its higher-dimensional analogues, and establish complete integrability of those maps.
\end{enumerate}
The first task presupposes a particular noncommutative leapfrog map, but its formula and domain are absent from the supplied record. The second task includes the construction of the maps themselves; it is not a theorem about a uniquely specified pre-existing noncommutative pentagram map.

As a commutative geometric reference, the planar pentagram construction takes a generic projective polygon $(v_i)$ to intersections of short diagonals, for example
\[
T(v)_i=\overline{v_{i-1}v_{i+1}}\cap\overline{v_iv_{i+2}},
\]
where the lines and intersections are defined. The source does not specify closed versus twisted polygons, quotient by projective transformations, dimension parameters, or which higher-dimensional constructions are intended. The formula is explanatory commutative notation, not a proposed noncommutative definition.

A rigorous noncommutative formulation requires a specified coefficient algebra or division ring, admissible elements and inverses, multiplication order, phase space, and notion of equivalence. It must also define ``complete integrability'' in that setting: ordinary Liouville integrability for a finite-dimensional Poisson realization, an appropriate noncommutative algebra of conserved quantities, a spectral/Lax formulation, or another stated criterion are different requirements. The supplied record chooses none of these, so no particular noncommutative map or integrability theorem is invented here.

The supplied summary reports established complete integrability of the commutative pentagram map and integrability of the leapfrog map in the KdV/Volterra family. It mentions noncommutative discrete systems and noncommutative/refined cluster connections, but reports no verified definitive published treatment resolving the two requested tasks. It labels this an open/unverified program, rather than a verified nonexistence result or proof of integrability.
\subsection{Short English statement}
Retakh asks first for integrability of a particular noncommutative leapfrog map, then for definitions and integrability proofs of noncommutative pentagram maps in the plane and higher dimensions. The supplied record omits the leapfrog formula and the noncommutative phase spaces and integrability convention, so the tasks remain a qualified research program.
\subsection{Sources}
\begin{itemize}
\item Local UnsolvedMath registry, permanent ID 2100409, exact source code catalogue entry, ``Open Problems in Integrable Systems --- Noncommutative integrable maps''; status open, classification OPEN-TRIAGE.
\item \emph{Open Problems, Questions, and Challenges in Finite-Dimensional Integrable Systems} (2018), Problem 4.9, attributed to V. Retakh, \url{https://arxiv.org/abs/1804.03737}. The record identifies extraction from public arXiv source TeX, accessed July 29, 2026, and says the source presented this as open.
\item The supplied research summary names Schwartz and Ovsienko--Schwartz--Tabachnikov for commutative pentagram integrability and cluster/Lax structures, and mentions higher-dimensional analogues ``by Glick etc.'' It also mentions the leapfrog KdV/Volterra connection and unspecified noncommutative/refined cluster works. No precise titles, dates, formulas, or theorem citations are supplied for these statements.
\item The summary says it could not verify a definitive published complete-integrability treatment of the noncommutative leapfrog map or definitions and proofs for the requested noncommutative pentagram maps. This is its supplied search assessment, not an independently verified literature conclusion.
\end{itemize}
\end{document}
